\clearpage
\SourceLeaf{374}
\pdfbookmark[2]{3. Gauss an Encke und Olbers}{tabula-nova-brief-003}
\begingroup
\makeatletter
\newcommand{\LetterThreeFirstNote}[1]{\begingroup\renewcommand{\thefootnote}{*)}\long\def\@makefntext##1{\noindent [*)\ ##1}\footnote{#1}\endgroup}
\makeatother
\begin{center}
[3.]\par\medskip
\textsc{Gauss} an [\textsc{Encke}, Göttingen 1815\LetterThreeFirstNote{Vgl. \textsc{Gauss} an \textsc{Olbers}, 29. Mai 1815, S. 350.]}.]\par\medskip
{\bfseries [Zur Berechnung der vorstehenden Tafel für die parabolische Bewegung.]}
\end{center}
\noindent\hspace*{\parindent}Auflösung der Gleichung
\[
\alpha\,\mathrm{tang}\,\tfrac12v+\beta\,\mathrm{tang}\,\tfrac12v^3=A,
\]
wo \(A\) durch seinen Logarithmus gegeben ist und
\[
\begin{aligned}
\alpha&=2\mathpunct{.}206265&&\text{u. genau}\quad\log\alpha=5{,}61545\ 51288\ 4044\\
\beta&=\tfrac23\mathpunct{.}206265&&\phantom{\text{u. genau}\quad}\log\beta=5{,}13833\ 38741\ 2078.
\end{aligned}
\]
Es sei \(V\) derjenige Winkel in runden \(20^{\mathrm{ern}}\) von Secunden, welcher \(v\) am nächsten kommt; man berechne
\[
\frac{\alpha\,\mathrm{tang}\,\tfrac12V}{\mathord{\text{\raisebox{-.35ex}{\(\scriptstyle:\)}}}A}
+\frac{\beta\,\mathrm{tang}\,\tfrac12V^3}{A}=m,
\]
welches von 1 sehr wenig verschieden sein wird. Sodann hat man
\[
v=V+(1-m)A\cos\tfrac12\left(\frac{v+V}{2}\right)^4;
\]
um diese Formel berechnen zu können, bedient man sich für \(v\) desjenigen Werths, den die \textsc{Barker}sche Tafel gibt\HistoricalNoteDoubleUnbracketed{Wenn es nöthig ist, kann man die Rechnung mit dem aus der Formel gefundenen verbesserten Werthe wiederholen.}, indem man mit \(\dfrac{75}{2\mathpunct{.}206265}\cdot A\) eingeht (oder mit dem Logarithm dieser Grösse). Hier ist
\[
\log\frac{75}{2\mathpunct{.}206265}=6{,}259\ 6061\,[-10].
\]
\par\medskip
\noindent\hspace*{2\parindent}Beispiel: \(\log A=6{,}166.\)
\par
\noindent\hspace*{\parindent}Hier gibt die \textsc{Barker}sche Tafel
\[
v=120^0\ 37'\ 18''{,}67,
\]
wir setzen also
\[
\begin{aligned}
V&=120\ \ 37\ \ 20\\
\tfrac14(V+v)&=\phantom1 60\ \ 18\ \ 39{,}6675
\end{aligned}
\]
\SourceLeaf{375}
\[
\begin{array}{l@{\;}l@{\;}l}
\mathrm{tang}\,\tfrac12V & .\,.\,.\,.\,.\,.\,.\,.\,.\,.\,.\,. & 0{,}24402\ 38664\\[3pt]
\dfrac{\alpha}{A} & .\,.\,.\,.\,.\,.\,.\,.\,.\,.\,.\,.\,.\,.\,.\,. & 9{,}44945\ 51288\\[3pt]\hline
&&9{,}69347\ 89952
\end{array}
\]
\[
\begin{array}{l@{\;}l@{\;}l}
\mathrm{tang}\,\tfrac12V^3 & .\,.\,.\,.\,.\,.\,.\,.\,.\,.\,. & 0{,}73207\ 15992\\[3pt]
\dfrac{\beta}{A} & .\,.\,.\,.\,.\,.\,.\,.\,.\,.\,.\,.\,.\,.\,.\,. & 8{,}97233\ 38741\\[3pt]\hline
&&9{,}70440\ 54733
\end{array}
\]
\[
\begin{array}{r@{\;}l}
\text{Zahl}&0{,}49371\ 80388\ 54\\
.\,.\,.\,.\,.\,.\,.\,.\,.&0{,}50629\ 71390\ 22\\\hline
m=&1{,}00001\ 51778\ 76
\end{array}
\]
\[
\begin{array}{l@{\;}l@{\;}l@{\quad}l}
1-m & .\,.\,.\,.\,.\,.\,.\,.\,.\,.\,.\,.\,.\,.\,.\,. &5{,}181\ 2109_n&\\[3pt]
\multicolumn{2}{r}{4\log\cos\dfrac{v+V}{4}=}&8{,}779\ 4436&\\[3pt]
A & .\,.\,.\,.\,.\,.\,.\,.\,.\,.\,.\,.\,.\,.\,.\,.\,.\,.\,.&6{,}166\ 0000&\\\cline{1-3}
&&0{,}126\ 6545_n&\text{Zahl }-1''{,}33861.
\end{array}
\]
Also
\[
v=120^0\ 37'\ 18''{,}66139.
\]
Die Wiederholung mit \(\dfrac{v+V}{4}=60^0\ 18'\ 34''{,}665\) gibt kein verändertes Resultat.
\par
\noindent\hspace*{\parindent}Um zu einem Logarithmen \(l\) die Zahl \(m\) zu finden, dient folgendes Verfahren. Es sei \(L\) der nächste Logarithm der Tafeln und \(M\) die dazu gehörige Zahl. Sodann hat man
\[
m=M+\frac{(l-L)\BareRoot{(mM)}}{k}\mathpunct{\text{\raisebox{.45ex}{,}}}
\]
wo \(k\) der Modulus der \textsc{Brigg}ischen Logarithmen und \(\log\dfrac1k=0{,}362\ 2157.\)
\par\nobreak
\noindent\hspace*{2\parindent}Beispiel:
\[
\begin{array}{r@{\;}l@{\;}l}
l&.\,.\,.\,.\,.&9{,}69347\ 90052\\
L&.\,.\,.\,.&9{,}69348\ 07203\\
&&0{,}.\,.\,.\,.\,.\,17151
\end{array}
\qquad M=0{,}49372
\]
\[
\begin{array}{l@{\;}l@{\;}l}
l-L&.\,.\,.\,.\,.\,.\,.\,.\,.\,.\,.\,.\,.\,.\,.\,.&4{,}234\ 2894_n\\[3pt]
\multicolumn{2}{r}{\tfrac12\log(mM)=}&9{,}693\ 4799\\[3pt]
\dfrac1k&.\,.\,.\,.\,.\,.\,.\,.\,.\,.\,.\,.\,.\,.\,.\,.\,.\,.\,.&0{,}362\ 2157\\[3pt]\hline
&&4{,}289\ 9850_n
\end{array}
\]
\[
\begin{array}{r@{\;}l}
\text{Zahl}&-0{,}.\,.\,.\,.\,.\,19497\ 8\\\hline
m=&0{,}49371\ 80502\ 2.
\end{array}
\]
\par
\noindent\hspace*{\parindent}Sie könnten nun den Anfang machen (um dies Beispiel mit benutzen zu \SourceLeaf{376}können) von \(\log A=6{,}166\), welches ungefähr in die Mitte der Tafel fällt, vorwärts zu gehen, indem Sie
\[
\begin{aligned}
\log A&=6{,}230\\
\log A&=6{,}294\\
\log A&=6{,}358\quad\text{etc.}
\end{aligned}
\]
setzen, während ich, so wie ich Zeit habe, rückwärts gehe, wozu Sie sich \textsc{Vegas} Thesaurus Logarithmorum von der Bibliothek nehmen können, da ich mein Exemplar hiebei selbst anwende. Die Gründe der Vorschriften werden Sie leicht selbst entwickeln können. Dies Blatt, wovon ich keine Abschrift behalten habe, erbitte ich mir bald zurück.\hfill G.
\par\medskip
\begin{center}\rule{.22\textwidth}{.4pt}\end{center}
\Needspace{8\baselineskip}
\pdfbookmark[3]{Gauss an Olbers, 16. Februar 1816}{tabula-nova-brief-003-olbers}
\begin{center}\textsc{Gauss} an \textsc{Olbers}, Göttingen, 16. Februar 1816.\end{center}
\noindent\hspace*{\parindent}.\ .\ .\ .\ .\quad Ich sehe aus Ihrem Briefe nicht bestimmt, ob \textsc{Burckhardt}s neue Tafel in der C[onn.] d[es] t[ems] abgedruckt oder bloss angezeigt ist. Im erstern Fall bitte ich mir ein paar auf einander folgende Glieder mitzutheilen. Ich selbst habe eine ganz ähnliche Tafel in müssigen Stunden angefangen und bereits zum grössern Theile vollendet. Wahrscheinlich ist aber doch noch einiger Unterschied in der Einrichtung. Ich habe die Tafel durchgehends auf 100000 Theile der Secunde berechnet, um die Hunderttheile (bis auf wenige Ausnahmen) immer gewiss zu geben; meine Argumente wachsen immer um 0,001, aber sie sind von den Logarithmen der Tage in \textsc{Lacailles} Tafel um eine Constante verschieden, die so gewählt ist, dass die absolute Zahl der Logarithmen für unendlich kleine Anomalien diesen selbst in Secunden gleich ist. Für die ersten \(11^0\) gebe ich die Zahlen selbst, wo dann das Interpoliren sehr leicht à vue geschehen kann, da die Differenz immer sehr wenig von 100 abweicht. Hier einige Proben\HistoricalNoteEight{\textsc{Gauss} gibt drei Proben, von denen hier nur die erste abgedruckt ist.]}:
\[
\renewcommand{\arraystretch}{.65}
\begin{array}{r@{\qquad}r@{\qquad}l}
\text{Argum.}&&\\[4pt]
8000&2^0\ 13'\ 18''{,}00&\\
&&99''{,}92\\
8100&14\ \ 57{,}92&\\
&&99{,}92\\
8200&16\ \ 37{,}84&\\
&&99{,}92\\
8300&18\ \ 17{,}76&\\
&&99{,}92.\\
8400&19\ \ 57{,}68&
\end{array}
\]
\par
\SourceLeaf{377}
\noindent\hspace*{\parindent}.\ .\ .\ .\ .\quad Die Tafel wird etwa bis \(167^0\) oder \(169^0\) fortgehen, und für den Schluss bis \(180^0\) wird wieder eine andere Einrichtung stattfinden, wobei man dieselbe Bequemlichkeit hat wie im Anfange, so dass die Tafel eine absolute Vollständigkeit hat und auch sonst für cubische Gleichungen mit Vortheil gebraucht werden kann. Bis jetzt ist die Tafel erst für die Hälfte der Argumente berechnet; durch eine bequeme Interpolationsmethode, die fast nichts als das Schreiben erfordert, wird das übrige ausgefüllt werden, so dass sie zusammen etwas über 4000 Glieder enthalten wird. Die Differenz \(348''{,}58\) ist das Maximum. Ich habe diese Tafel schon im vorigen Jahre bei Ihrem Cometen mit Vortheil angewandt.\quad .\ .\ .\ .\ .
\par\medskip
\begin{center}\rule{.22\textwidth}{.4pt}\end{center}
\endgroup
